\section{Fase-Fase Kompilasi: Analisis dan Sintesis}
\label{sec:ch03-fase}

Kompilator modern biasanya dibagi menjadi dua bagian besar: analisis
(front-end) dan sintesis (back-end). Bagian analisis memahami struktur
program sumber; bagian sintesis menghasilkan program target. Batas
keduanya adalah representasi antara (\eng{intermediate
representation}, IR).

\begin{figure}[H]
\centering
\begin{tikzpicture}[node distance=6mm and 7mm]
  \node[phase, text width=2.3cm] (src) {Program\\sumber};
  \node[phase, text width=2.3cm, right=of src] (lex) {Analisis\\leksikal};
  \node[phase, text width=2.3cm, right=of lex] (syn) {Analisis\\sintaksis};
  \node[phase, text width=2.3cm, right=of syn] (sem) {Analisis\\semantik};
  \node[data, text width=1.6cm, right=of sem] (ir) {IR};
  \node[phase, text width=2.3cm, below=of ir] (opt) {Optimasi};
  \node[phase, text width=2.6cm, left=of opt] (gen) {Pembangkitan\\kode};
  \node[phase, text width=2.3cm, left=of gen] (tgt) {Program\\target};
  \draw[arrow] (src) -- (lex);
  \draw[arrow] (lex) -- (syn);
  \draw[arrow] (syn) -- (sem);
  \draw[arrow] (sem) -- (ir);
  \draw[arrow] (ir) -- (opt);
  \draw[arrow] (opt) -- (gen);
  \draw[arrow] (gen) -- (tgt);
\end{tikzpicture}
\caption{Pipeline kompilasi dari program sumber ke program target}
\label{fig:ch03-pipeline}
\end{figure}

\subsection{Analisis Leksikal}

Analisis leksikal membaca karakter demi karakter dan mengelompokkannya
menjadi token. Token memiliki jenis dan atribut, misalnya
\code{INT\_LIT} dengan nilai 42. Spasi, baris baru, dan komentar
dibuang. Keluaran fase ini adalah aliran token yang siap diproses
parser.

\subsection{Analisis Sintaksis}

Analisis sintaksis memeriksa apakah aliran token membentuk struktur
yang sesuai dengan tata bahasa. Hasilnya adalah pohon urai atau AST.
Parser juga dapat membangun tabel simbol awal dan melakukan pemulihan
kesalahan sintaksis.

\subsection{Analisis Semantik}

Analisis semantik memeriksa makna program: apakah identifier sudah
dideklarasikan, apakah tipe operand cocok, apakah fungsi dipanggil
dengan jumlah argumen yang benar, dan apakah kontrol alur valid.
Keluaran fase ini adalah AST yang telah dianotasi dengan informasi
tipe dan referensi simbol.

\subsection{Sintesis}

Setelah IR terbentuk, fase sintesis melakukan optimasi dan
pembangkitan kode. Optimasi memperbaiki IR tanpa mengubah makna
program; pembangkitan kode memilih instruksi target,
mengalokasikan register, dan menghasilkan program assembly atau
kode mesin. Teknik optimasi tidak dirinci pada bab ini karena fokus
mata kuliah adalah \eng{front-end}.

\begin{table}[H]
\centering
\caption{Masukan dan keluaran setiap fase}
\label{tab:ch03-fase}
\begin{tabularx}{\textwidth}{@{}lXX@{}}
\toprule
\textbf{Fase} & \textbf{Masukan} & \textbf{Keluaran} \\
\midrule
Leksikal & Karakter sumber & Aliran token \\
Sintaksis & Aliran token & Pohon urai / AST \\
Semantik & AST dan tabel simbol & AST teranotasi \\
IR & AST teranotasi & Kode antara \\
Optimasi & Kode antara & Kode antara yang lebih efisien \\
Pembangkitan kode & Kode antara & Assembly atau kode mesin \\
\bottomrule
\end{tabularx}
\end{table}

\begin{contoh}[Perjalanan program sederhana]
Untuk program
\begin{lstlisting}[language=C]
int main() {
    int a = 5;
    int b = 7;
    print(a + b);
    return 0;
}
\end{lstlisting}
fase leksikal menghasilkan token \code{int}, \code{main}, \code{(},
\code{)}, \code{\{}, dan seterusnya. Fase sintaksis membentuk pohon
deklarasi fungsi dan pernyataan. Fase semantik mencatat bahwa
\code{a} dan \code{b} bertipe \code{int}, serta bahwa \code{print}
menerima satu argumen. Fase IR dapat menghasilkan instruksi
tiga-alamat seperti \code{t1 = a + b} dan \code{print t1}.
\end{contoh}
